\documentclass[11pt,letterpaper]{article}
\usepackage[utf8]{inputenc}
\usepackage[T1]{fontenc}
\usepackage{lmodern}
\usepackage[english]{babel}
\usepackage[margin=0.85in]{geometry}
\usepackage{enumitem}
\usepackage{titlesec}
\usepackage{tabularx}
\usepackage[hidelinks]{hyperref}

\hypersetup{pdftitle={Anicet Nzeukoc -- CV}, pdfauthor={Anicet Nzeukoc}}
\pagestyle{empty}
\setlength{\parindent}{0pt}

% Section titles: uppercase, bold, full-width rule underneath
\titleformat{\section}{\normalsize\bfseries\uppercase}{}{0pt}{}[\vspace{-3pt}\hrule height 0.4pt]
\titlespacing*{\section}{0pt}{12pt}{6pt}

% Section content is indented; bullets (centred dot) hang in the indentation
\newenvironment{block}{\begin{list}{}{\setlength{\leftmargin}{1.6em}\setlength{\topsep}{0pt}%
  \setlength{\partopsep}{0pt}\setlength{\itemsep}{0pt}\setlength{\parsep}{0pt}}\item[]}{\end{list}}
\setlist[itemize]{label=\textperiodcentered, leftmargin=0pt, labelsep=1.1em, itemsep=1pt, topsep=2pt, parsep=0pt}

% \entry{Title}{Dates}: bold title, italic dates on the right
\newcommand{\entry}[2]{\textbf{#1}\hfill\textit{#2}\par}
% \job{Organisation}{Dates}{Role}{Place}
\newcommand{\job}[4]{\textbf{#1}\hfill #2\par\textit{#3}\hfill\textit{#4}\par}
\newcommand{\gap}{\par\vspace{0.6em}}

\begin{document}

% -------------------- HEADER --------------------
\begin{center}
{\LARGE\bfseries ANICET NZEUKOC}
\end{center}
\vspace{1.4em}

\begin{minipage}[t]{0.52\textwidth}
University of Montreal (UdeM)\\
Department of Economics (CIREQ)\\
3150 Jean-Brillant Street\\
Montreal (Quebec, Canada), H3C 3J7\\
Office: C6070-10
\end{minipage}%
\begin{minipage}[t]{0.48\textwidth}
Email: \href{mailto:nzeukoca@gmail.com}{nzeukoca@gmail.com}\\
Website: \href{https://anicetnzeu.github.io}{anicetnzeu.github.io}\\
GitHub: \href{https://github.com/AnicetNzeu}{github.com/AnicetNzeu}\\
LinkedIn: \href{https://www.linkedin.com/in/anicet-nzeukoc-660000156}{in/anicet-nzeukoc-660000156}
\end{minipage}

\vspace{1.6em}
\rule{0.5\textwidth}{0.4pt}

% -------------------- SUMMARY --------------------
\section*{Summary}
\begin{block}
I am a PhD candidate in Economics at the University of Montreal and a member of CIREQ, supervised by
Professors Marine Carrasco and René Garcia. My research fields are econometrics and financial
econometrics: functional data analysis, structural breaks, option pricing and risk-neutral densities,
and machine learning methods for structural models. I develop instrumental-variable and change-point
methods to recover and monitor the probability distributions that option markets assign to future asset
prices. I bring 5+ years of experience in quantitative research across Cameroon and Canada, and I work
in French and English.
\end{block}

% -------------------- EDUCATION --------------------
\section*{Education}
\begin{block}
\entry{University of Montreal (UdeM), CIREQ}{2022--}
PhD candidate in Economics, Canada. Supervisors: Prof.\ Marine Carrasco and Prof.\ René Garcia
\gap
\entry{University of Yaoundé II}{2021--2022}
MSc in Economics, Cameroon
\gap
\entry{Sub-regional Institute of Statistics and Applied Economics (ISSEA)}{2018--2021}
Engineer in Statistics and Economics (ISE), Cameroon
\gap
\entry{University of Dschang}{2016--2018}
MSc in Applied Mathematics, Cameroon
\gap
\entry{University of Dschang}{2013--2016}
BSc (Licence) in Mathematics and Computer Science, Cameroon
\end{block}

% -------------------- EXPERIENCE --------------------
\section*{Professional and Research Experience}
\begin{block}
\job{Télé-université (TELUQ)}{April 2026 -- August 2026}{Research Assistant (part-time)}{Montréal}
\begin{itemize}
\item Governance, directors' and officers' insurance and executive compensation.
\item Built a firm-year panel on more than 220 Canadian public firms (2013--2024) with an automated
      extraction pipeline in Python.
\item Applied difference-in-differences and event-study methods.
\end{itemize}
\gap
\job{HEC Montréal}{December 2022 -- August 2025}{Research Assistant, with Prof.\ Pierre-Olivier Pineau}{Montréal}
\begin{itemize}
\item Stochastic frontier analysis of energy efficiency across the ten Canadian provinces.
\item Drafted ``Measuring Provincial Energy Efficiency in Canada: A Stochastic Frontier Analysis''.
\end{itemize}
\gap
\job{Scikoop Inc.\ / IVADO}{January 2023 -- June 2023}{Research Consultant, IVADO Scientist in Residence program}{Montréal}
\begin{itemize}
\item Statistical tools for equity, diversity and inclusion in public-health and social-science research.
\end{itemize}
\gap
\job{ISSEA}{October 2021 -- July 2022}{Research Assistant, with Dr.\ Ignace Kamga}{Yaoundé}
\begin{itemize}
\item Macroeconomic forecasting for the CEMAC countries.
\item Organised the weekly ``Research Wednesday'' seminar.
\end{itemize}
\gap
\entry{Teaching}{2024--2025}
Lecturer, Econometrics I (UdeM, ECN1260); Fall 2024\\
Lecturer, Principles of Economics (UdeM, ECN1000); Fall 2025
\gap
\entry{Teaching assistant}{2021--}
Economic Data Analysis (UdeM, ECN1160); Fall 2025 and Fall 2026, for Prof.\ Marine Carrasco\\
Advanced Econometrics Theory (UdeM, ECN7060, PhD); Winter 2026, for Prof.\ Marine Carrasco\\
Numerical Methods in Economics (UdeM, ECN6338); Winter and Fall 2025, for Prof.\ William McCausland\\
Techniques of Economic Analysis II (UdeM, ECN1075); Summer 2025, for Prof.\ Eugène Deta'a\\
Econometrics I (UdeM, ECN1260); Winter 2025, for Prof.\ Amal Serre\\
Applied Econometrics (UdeM, ECN2260); Winter 2024, for Prof.\ William Arbour\\
Microeconomic Theory I (UdeM, ECN2040); Winter and Fall 2023, for Prof.\ Andriana Bellou\\
Microeconomics (ISSEA, engineering-economist students); 2021--2022, for Dr.\ Ignace Kamga
\gap
\entry{Research in progress}{}
``Instrumental Variable Estimation of Risk-Neutral Density using Functional Linear Models,''
with Marine Carrasco and René Garcia.\\
``Deep Surrogates in Option Pricing Models for Risk-Neutral Density Estimation,'' with Rostand Tchouakam.\\
``Detecting and Locating Structural Changes in a Sequence of Estimated Risk-Neutral Densities,''
work in progress.
\end{block}

% -------------------- AWARDS --------------------
\section*{Awards and Fellowships}
\begin{block}
Fonds de recherche du Québec -- Société et culture (FRQSC), Doctoral Research Fellowship (B2), 2026--2027\\
First Prize, Poster Competition, 90th ACFAS Conference, Université de Montréal, 2023\\
Doctoral Admission Fellowship, Department of Economics, UdeM, 2022--\\
Tuition Fee Exemption Scholarship, University of Montreal, 2022--\\
Government of Cameroon Admission Scholarship, ISSEA, 2018--2021
\end{block}

% -------------------- AFFILIATIONS --------------------
\section*{Affiliations}
\begin{block}
Centre interuniversitaire de recherche en économie quantitative (CIREQ), 2023--\\
Society for Financial Econometrics (SoFiE), 2024--
\end{block}

% -------------------- TECHNICAL STRENGTHS --------------------
\section*{Technical Strengths}
\begin{block}
\begin{tabularx}{\linewidth}{@{}l X@{}}
\textbf{Programming} & Python (NumPy, Pandas, PyTorch, QuantLib, scikit-learn), R, MATLAB, Stata, SQL, Git \\
\textbf{Data}        & OptionMetrics Ivy DB, large-scale time series, parallel computing, reproducible pipelines \\
\textbf{Office}      & \LaTeX, Excel (advanced), PowerPoint \\
\end{tabularx}
\end{block}

% -------------------- SERVICE AND TALKS --------------------
\section*{Volunteering Activities and Talks}
\begin{block}
\entry{Chair, Quebec PhD Workshop in Economics, Statistics and Finance (QPESF)}{2022--}
Organiser and moderator of a weekly doctoral seminar in economics, statistics and finance across Quebec
universities.
\gap
\entry{Co-organiser, CIREQ Doctoral Conference on Economics and Finance}{April 2026}
Organising committee with Prof.\ Marine Carrasco (UdeM), Jackson Montelpare (McGill) and Zahraa Arkahdan
(Concordia).
\gap
\entry{Volunteer, Cap Campus (Université de Montréal)}{2024--}
Workshops in applied mathematics, finance and statistics for high-school and college students.
\gap
\entry{Talks}{2023--2026}
41st Annual Meeting of the Canadian Econometrics Study Group (CESG), Simon Fraser University,
Vancouver; October 2026 (poster)\\
65th Annual Congress of the Société canadienne de science économique (SCSE), Québec City; May 2026\\
CIREQ Doctoral Conference on Econometrics and Finance, Université de Montréal; April 2025\\
Doctoral seminar, Department of Economics, Université de Montréal; May 2025 and December 2025\\
Inclusive Society Colloquium, Université Laval, Québec; June 2023\\
ACFAS Conference, 90th edition, Université de Montréal; May 2023 (poster)
\end{block}

% -------------------- LANGUAGES --------------------
\section*{Language and Other Interests}
\begin{block}
French (native), English (fluent), football, reading, writing
\end{block}

% -------------------- ABSTRACTS --------------------
\section*{Abstract of Some Selected Work in Progress}
\begin{block}
\textbf{Instrumental Variable Estimation of Risk-Neutral Density using Functional Linear Models}\\
With Marine Carrasco and René Garcia\\
\textit{Nonparametric estimation of the risk-neutral density from a time series of European option prices.
Pooling contracts across dates enlarges the sample but makes the strike endogenous, because the exchange
sets strikes in response to the same shocks that move prices. The paper develops an instrumental-variable
sieve estimator in Hermite basis functions, using lagged values of the index and of the VIX and SKEW
indices as instruments, together with an estimator-free confidence band for the entire density that
accommodates serially dependent pricing errors. The Monte Carlo design is the estimated model written in
reduced form, so the experiment reproduces the measured degree of endogeneity. Applied to S\&P 500
options, OptionMetrics, 1996--2022.}

\vspace{0.4em}
\hspace*{4em}\textbf{JEL Classification:} C14, C26, C58, G13\\
\hspace*{4em}\textbf{Keywords:} Risk-Neutral Density, Endogeneity, Instrumental Variables, Sieve Estimation
\gap\gap
\textbf{Deep Surrogates in Option Pricing Models for Risk-Neutral Density Estimation}\\
With Rostand Tchouakam
\gap\gap
\textbf{Detecting and Locating Structural Changes in a Sequence of Estimated Risk-Neutral Densities}\\
Sole-authored
\end{block}

\end{document}
